% ===========================================================================
%  Chapitre 4 — Primitives et calcul intégral
%  PE 2026, composante ANALYSE
% ===========================================================================
\bmtchapitre{Primitives et calcul intégral}
  {Calculer une primitive et une intégrale ; intégrer par parties et par
   changement de variable affine ; calculer une aire et une valeur moyenne.}
  {Analyse \textbullet\ environ 12 heures}

Jusqu'ici, une fonction étant donnée, vous calculiez sa dérivée. Ce chapitre
prend le problème à l'envers : une dérivée étant donnée, retrouver la fonction
dont elle provient. L'opération n'a rien d'automatique — dériver est une
mécanique, primitiver est une reconnaissance.

Et il se produit alors quelque chose d'inattendu. Cette question purement
algébrique donne la réponse à une question de géométrie que personne
n'aurait crue liée : quelle est l'aire sous une courbe. Cette même idée
permettra, plus tard dans l'année, de calculer une valeur moyenne — celle
d'un chiffre d'affaires, d'une production, d'un coût — sur toute une période.

% ---------------------------------------------------------------------------
\begin{revision}
Uniquement des dérivées — c'est le sens que ce chapitre va renverser.

\begin{enumerate}[label=\textbf{\arabic*.}, leftmargin=*, itemsep=0.35em]
  \item Dériver : $x \mapsto x^{5}$, $x \mapsto \sqrt{x}$,
        $x \mapsto \dfrac{1}{x^{2}}$, $x \mapsto \sin(2x)$.
  \item Dériver $x \mapsto (3x+1)^{4}$, puis $x \mapsto \sqrt{2x-5}$.
  \item Rappeler la formule de dérivation d'un produit.
  \item Rappeler la formule $\left(u^{n}\right)'$ pour $n$ entier.
  \item Quelle est l'aire d'un rectangle de largeur $b-a$ et de hauteur $h$ ?
  \item Rappeler la définition d'une primitive.
\end{enumerate}
\end{revision}

\begin{automatismes}
Sans calculatrice, sans brouillon.

\begin{enumerate}[label=\textbf{\alph*)}, leftmargin=*, itemsep=0.25em]
  \item Dérivée de $x \mapsto \dfrac{x^{4}}{4}$
  \item Dérivée de $x \mapsto 2\sqrt{x}$
  \item Dérivée de $x \mapsto -\dfrac{1}{x}$
  \item Dérivée de $x \mapsto -\dfrac{\cos(3x)}{3}$
  \item Dérivée de $x \mapsto \dfrac{(2x+1)^{5}}{10}$
  \item Une fonction dont la dérivée est $x \mapsto x$
  \item Une fonction dont la dérivée est $x \mapsto \cos x$
  \item Une fonction dont la dérivée est $x \mapsto 3x^{2}$
\end{enumerate}
\end{automatismes}

\begin{corrige}[automatismes]
a) $x^{3}$ \quad b) $\dfrac{1}{\sqrt{x}}$ \quad c) $\dfrac{1}{x^{2}}$ \quad
d) $\sin(3x)$ \quad e) $(2x+1)^{4}$ \quad f) $\dfrac{x^{2}}{2}$ \quad
g) $\sin x$ \quad h) $x^{3}$.
\end{corrige}

% ===========================================================================
\section{Les primitives}

\begin{definition}[primitive]
Soit $f$ une fonction définie sur un intervalle $I$. On appelle
\textbf{primitive} de $f$ sur $I$ toute fonction $F$ dérivable sur $I$ telle
que
\[
  F'(x) = f(x) \qquad \text{pour tout } x \in I .
\]
\end{definition}

\begin{theoreme}[existence et forme générale]
Toute fonction continue sur un intervalle $I$ admet des primitives sur $I$.
Si $F$ en est une, les primitives de $f$ sur $I$ sont exactement les
fonctions $x \mapsto F(x)+k$, où $k$ décrit $\R$.
\end{theoreme}

\begin{demonstration}
Si $G$ est une autre primitive, $H = G-F$ vérifie $H' = 0$ sur $I$. Une
fonction de dérivée nulle sur un intervalle y est constante — conséquence du
théorème des accroissements finis, appliqué entre deux points quelconques $a$
et $b$ de $I$ : $H(b)-H(a) = (b-a)H'(c) = 0$. Donc $G = F+k$.
\end{demonstration}

\subsection{Le tableau des primitives usuelles}

\begin{center}
\renewcommand{\arraystretch}{1.6}
\begin{tabular}{@{}p{4.3cm} p{4.3cm} p{3.6cm}@{}}
\toprule
\textsf{\footnotesize\bfseries FONCTION $f$} &
\textsf{\footnotesize\bfseries UNE PRIMITIVE $F$} &
\textsf{\footnotesize\bfseries SUR} \\
\midrule
$x \mapsto a$ (constante)        & $x \mapsto ax$                       & $\R$ \\
$x \mapsto x^{n}$, $n \in \N$    & $x \mapsto \dfrac{x^{n+1}}{n+1}$     & $\R$ \\
$x \mapsto \dfrac{1}{x^{2}}$     & $x \mapsto -\dfrac{1}{x}$            & tout intervalle de $\Rs$ \\
$x \mapsto \dfrac{1}{\sqrt{x}}$  & $x \mapsto 2\sqrt{x}$                & $\Rps$ \\
\midrule
$u'u^{n}$, $n \neq -1$           & $\dfrac{u^{n+1}}{n+1}$               & où $u$ est dérivable \\
$\dfrac{u'}{\sqrt{u}}$           & $2\sqrt{u}$                          & où $u > 0$ \\
$u' \times (v'\circ u)$          & $v \circ u$                          & \\
\bottomrule
\end{tabular}
\end{center}

\bmtmarge{La seconde moitié du tableau est la plus utile : reconnaître un
$u'u^{n}$ dans une expression compliquée est la compétence centrale du
chapitre.}

\begin{exemple}[reconnaître la forme $u'u^{n}$]
Cherchons une primitive de $f(x) = 3x^{2}\left(x^{3}+1\right)^{4}$ sur $\R$.

On pose $u(x) = x^{3}+1$, donc $u'(x) = 3x^{2}$. L'expression est exactement
$u'u^{4}$, dont une primitive est $\dfrac{u^{5}}{5}$. Donc
\[
  F(x) = \frac{\left(x^{3}+1\right)^{5}}{5}
\]
convient, et l'on vérifie en dérivant.
\end{exemple}

\begin{propriete}[primitive prenant une valeur donnée en un point]
Soit $f$ continue sur $I$, $x_{0} \in I$ et $y_{0}$ un réel. Il existe une
\emph{unique} primitive $F$ de $f$ sur $I$ telle que $F(x_{0}) = y_{0}$.
\end{propriete}

\begin{entrainement}[primitives]
\exo[\niveau{1}] Déterminer une primitive de $x \mapsto 4x^{3}-2x+5$ sur $\R$,
puis de $x \mapsto \dfrac{3}{x^{2}}$ sur $\Rps$.

\exo[\niveau{2}] Déterminer une primitive de
$x \mapsto \dfrac{2x}{\sqrt{x^{2}+1}}$ sur $\R$.

\exo[\niveau{3}] Déterminer la primitive $F$ de
$x \mapsto \dfrac{x}{\left(x^{2}+1\right)^{2}}$ sur $\R$ telle que $F(0) = 1$.
\end{entrainement}

\begin{corrige}
\textbf{1.} $x^{4}-x^{2}+5x$ ; \ $-\dfrac{3}{x}$.

\textbf{2.} C'est $\dfrac{u'}{\sqrt{u}}$ avec $u = x^{2}+1$ : une primitive est
$2\sqrt{x^{2}+1}$.

\textbf{3.} L'expression est $\dfrac{1}{2}u'u^{-2}$ avec $u = x^{2}+1$, de
primitive $-\dfrac{1}{2u}$. La condition $F(0)=1$ donne $k = \tfrac32$ :
$F(x) = \dfrac{-1}{2(x^{2}+1)} + \dfrac{3}{2}$.
\end{corrige}

% ===========================================================================
\section{L'intégrale}

\begin{definition}[intégrale d'une fonction continue]
Soit $f$ continue sur un intervalle $I$, et $a$, $b$ deux éléments de $I$. Si
$F$ est une primitive quelconque de $f$ sur $I$, on pose
\[
  \int_{a}^{b} f(x)\dx = F(b)-F(a),
\]
quantité que l'on note aussi $\bigl[F(x)\bigr]_{a}^{b}$.
\end{definition}

\begin{propriete}[règles de calcul]
Pour $f$ et $g$ continues sur $I$, $a$, $b$, $c$ dans $I$ et $\lambda$ réel :
\begin{align*}
  \int_{a}^{a} f(x)\dx &= 0, &
  \int_{b}^{a} f(x)\dx &= -\int_{a}^{b} f(x)\dx, \\[0.3em]
  \int_{a}^{b} \bigl[f(x)+g(x)\bigr]\dx &= \int_{a}^{b} f + \int_{a}^{b} g, &
  \int_{a}^{b} \lambda f(x)\dx &= \lambda \int_{a}^{b} f(x)\dx,
\end{align*}
et la \textbf{relation de Chasles} :
$\displaystyle\int_{a}^{b} f(x)\dx = \int_{a}^{c} f(x)\dx + \int_{c}^{b} f(x)\dx$.
\end{propriete}

\begin{propriete}[positivité et ordre]
Si $a \leq b$ :
\begin{itemize}[leftmargin=1.4em, itemsep=0.18em]
  \item $f \geq 0$ sur $\intervalle{a}{b}$ entraîne $\displaystyle\int_{a}^{b} f(x)\dx \geq 0$ ;
  \item $f \leq g$ sur $\intervalle{a}{b}$ entraîne
        $\displaystyle\int_{a}^{b} f(x)\dx \leq \int_{a}^{b} g(x)\dx$ ;
  \item si $m \leq f(x) \leq M$ sur $\intervalle{a}{b}$, alors
        $m(b-a) \leq \displaystyle\int_{a}^{b} f(x)\dx \leq M(b-a)$.
\end{itemize}
\end{propriete}

% ===========================================================================
\section{Deux techniques de calcul}

\subsection{L'intégration par parties}

\begin{theoreme}[intégration par parties]
Soient $u$ et $v$ deux fonctions dérivables sur $\intervalle{a}{b}$, de
dérivées continues. Alors
\[
  \int_{a}^{b} u(x)v'(x)\dx
  = \bigl[u(x)v(x)\bigr]_{a}^{b} - \int_{a}^{b} u'(x)v(x)\dx .
\]
\end{theoreme}

\begin{demonstration}
La formule de dérivation d'un produit donne $(uv)' = u'v + uv'$. On intègre
cette égalité entre $a$ et $b$ :
$\bigl[u(x)v(x)\bigr]_{a}^{b} = \int_{a}^{b} u'(x)v(x)\dx + \int_{a}^{b} u(x)v'(x)\dx$,
puis on isole la seconde intégrale.
\end{demonstration}

\begin{methode}{choisir $u$ et $v'$}
On prend pour $u$ la fonction que la dérivation \emph{simplifie}, et pour
$v'$ celle que l'on sait primitiver sans peine. Un polynôme multiplié par un
cosinus, un sinus ou une exponentielle : le polynôme est $u$, car sa dérivée
baisse d'un degré et finit par disparaître.
\end{methode}

\begin{exemple}[une intégration par parties menée en entier]
Calculons $\displaystyle\int_{0}^{\pi} x\sin x \dx$.

On pose $u(x) = x$ et $v'(x) = \sin x$, d'où $u'(x) = 1$ et
$v(x) = -\cos x$. La formule donne
\[
  \int_{0}^{\pi} x\sin x \dx
  = \bigl[-x\cos x\bigr]_{0}^{\pi} + \bigl[\sin x\bigr]_{0}^{\pi}
  = \pi + 0 = \pi .
\]
\end{exemple}

\subsection{Le changement de variable affine}

\begin{propriete}[changement de variable affine]
Soit $f$ continue et $F$ une primitive de $f$. Pour $\alpha \neq 0$,
\[
  \int_{a}^{b} f(\alpha x + \beta)\dx
  = \frac{1}{\alpha}\Bigl[F(\alpha x + \beta)\Bigr]_{a}^{b} .
\]
\end{propriete}

\begin{entrainement}[calcul d'intégrales]
\exo[\niveau{1}] Calculer $\displaystyle\int_{1}^{2}\left(3x^{2}-4x+1\right)\dx$.

\exo[\niveau{2}] Calculer $\displaystyle\int_{0}^{\frac{\pi}{2}} x\cos x \dx$
par une intégration par parties.
\end{entrainement}

\begin{corrige}
\textbf{1.} $\bigl[x^{3}-2x^{2}+x\bigr]_{1}^{2} = 2-0 = 2$.

\textbf{2.} $u=x$, $v'=\cos x$ : $\bigl[x\sin x\bigr]_{0}^{\pi/2}
- \int_{0}^{\pi/2}\sin x\dx = \dfrac{\pi}{2} - 1$.
\end{corrige}

% ===========================================================================
\section{Aires et valeur moyenne}

\begin{theoreme}[interprétation d'une intégrale comme une aire]
Soit $f$ continue et \emph{positive} sur $\intervalle{a}{b}$. Alors
$\displaystyle\int_{a}^{b} f(x)\dx$ est l'aire, exprimée en unités d'aire, du
domaine plan défini par $a \leq x \leq b$ et $0 \leq y \leq f(x)$.
\end{theoreme}

\begin{visualisation}[l'aire sous la courbe]
\begin{center}
\begin{tikzpicture}
\begin{axis}[
  width = 10.6cm, height = 5.4cm,
  axis lines = middle, xlabel = {$x$},
  xmin = -0.5, xmax = 5.4, ymin = -0.4, ymax = 3.4,
  xtick = {1,4}, xticklabels = {$a$,$b$}, ytick = \empty,
  axis line style = {bmtGris},
  tick label style = {font=\footnotesize, color=bmtGris}, clip = true,
]
  \addplot[bmtIndigoPale, draw=none, fill=bmtIndigo, fill opacity=0.16,
           domain=1:4, samples=100] {0.35*x^2 - 1.6*x + 3.1} \closedcycle;
  \addplot[bmtIndigo, line width=1.3pt, domain=0.2:5.2, samples=140]
    {0.35*x^2 - 1.6*x + 3.1};
  \draw[bmtGris, dotted] (axis cs:1,0) -- (axis cs:1,1.85);
  \draw[bmtGris, dotted] (axis cs:4,0) -- (axis cs:4,2.3);
  \node[bmtIndigo, font=\footnotesize] at (axis cs:2.5,0.72)
    {$\displaystyle\int_{a}^{b} f(x)\,dx$};
\end{axis}
\end{tikzpicture}
\end{center}

La partie teintée est bornée en bas par l'axe des abscisses, en haut par la
courbe. Si la fonction change de signe, découpez avec la relation de Chasles
et prenez la valeur absolue sur chaque morceau.
\end{visualisation}

\begin{definition}[valeur moyenne]
La \textbf{valeur moyenne} de $f$ sur $\intervalle{a}{b}$, avec $a<b$, est le
réel
\[
  \mu = \frac{1}{b-a}\int_{a}^{b} f(x)\dx .
\]
\end{definition}

\begin{remarque}
Géométriquement, $\mu$ est la hauteur du rectangle de base $\intervalle{a}{b}$
qui aurait la même aire que le domaine sous la courbe. En économie, si $f(t)$
représente un chiffre d'affaires instantané ou une production au temps $t$,
$\mu$ en donne la valeur moyenne sur toute la période — un indicateur que les
données brutes, dispersées mois par mois, ne montrent pas directement.
\end{remarque}

\begin{entrainement}[aires et valeurs moyennes]
\exo[\niveau{1}] Calculer l'aire du domaine délimité par la courbe de
$f(x) = x^{2}$, l'axe des abscisses et les droites $x=0$ et $x=3$.

\exo[\niveau{2}] Calculer la valeur moyenne de $x \mapsto x^{2}$ sur
$\intervalle{0}{3}$, et vérifier qu'elle est comprise entre le minimum et le
maximum de la fonction sur cet intervalle.
\end{entrainement}

\begin{corrige}
\textbf{1.} $\left[\dfrac{x^{3}}{3}\right]_{0}^{3} = 9$ unités d'aire.

\textbf{2.} $\mu = \dfrac{1}{3}\times 9 = 3$, comprise entre $0$ et $9$.
\end{corrige}

% ===========================================================================
% ===========================================================================
\begin{annales}
Les énoncés rassemblés ici sont repris de sujets réellement posés ; leur
provenance est indiquée à chaque fois. Aucun ne fait appel au logarithme ni à
l'exponentielle.

\exoann{Baccalauréat, Côte d'Ivoire, séries C et D, session 2023 — exercice 2}
Soit $f$ la fonction définie sur $\R$ par $f(x) = \cos x - x\sin x$. Déterminer
une primitive de $f$ sur $\R$.

\exoann{Baccalauréat, Côte d'Ivoire, séries C et D, session 2023 — exercice 2, suite}
En déduire la valeur de
$\displaystyle\int_{0}^{\frac{\pi}{2}}\left(\cos x - x\sin x\right)\dx$.
\end{annales}

\begin{corrige}[annales]
\textbf{1.} Le produit $x\cos x$ se dérive en $\cos x - x\sin x$, qui est
exactement $f(x)$. Les primitives de $f$ sur $\R$ sont donc les fonctions
$F(x) = x\cos x + k$, $k\in\R$.

\textbf{2.} Avec $F(x) = x\cos x$,
$\displaystyle\int_{0}^{\pi/2}\left(\cos x - x\sin x\right)\dx
  = \left[x\cos x\right]_{0}^{\pi/2} = 0$.
\end{corrige}

\begin{jecherche}[le silo à grains de Marovoay]
La plaine de Marovoay possède un silo obtenu en faisant tourner, autour de
l'axe des abscisses, le domaine sous la courbe de
$f(x) = \sqrt{9-x^{2}}$, $x \in \intervalle{0}{3}$, les longueurs étant
exprimées en mètres. On admet que le volume ainsi engendré vaut
$V = \pi\displaystyle\int_{0}^{3}\left[f(x)\right]^{2}\dx$.

\begin{enumerate}[label=\textbf{\arabic*.}, leftmargin=*, itemsep=0.3em]
  \item Vérifier que $f$ est bien définie et positive sur $\intervalle{0}{3}$.
  \item Calculer le volume $V$ du silo, en mètres cubes.
  \item Le prix de revient de la construction est proportionnel au volume, au
        taux de $45\,000$ Ar par mètre cube. Estimer ce coût, en prenant
        $\pi \approx 3{,}14$.
  \item La récolte quotidienne livrée au silo, en tonnes, est modélisée sur
        une semaine de sept jours par $g(t) = 2+0{,}5t$, $t \in \intervalle{0}{7}$.
        Calculer la valeur moyenne de $g$ sur cette semaine, et interpréter.
\end{enumerate}
\end{jecherche}

\begin{corrige}[le silo à grains de Marovoay]
\textbf{1.} $9-x^{2} \geq 0$ équivaut à $x^{2} \leq 9$, vrai sur
$\intervalle{0}{3}$ : $f$ y est définie et positive.

\textbf{2.} $V = \pi\displaystyle\int_{0}^{3}\left(9-x^{2}\right)\dx
= \pi\left[9x-\frac{x^{3}}{3}\right]_{0}^{3} = \pi(27-9) = 18\pi \approx 56{,}5$
m$^{3}$.

\textbf{3.} Coût $\approx 56{,}5 \times 45\,000 \approx 2\,542\,500$ Ar.

\textbf{4.} $\mu = \dfrac{1}{7}\displaystyle\int_{0}^{7}(2+0{,}5t)\dt
= \dfrac{1}{7}\left[2t+0{,}25t^{2}\right]_{0}^{7}
= \dfrac{1}{7}(14+12{,}25) \approx 3{,}75$ tonnes par jour : c'est la
livraison quotidienne moyenne sur la semaine, alors même que la livraison
réelle croît chaque jour.
\end{corrige}

% ===========================================================================
\begin{jemeteste}
Répondre par vrai ou faux, en justifiant en une ligne.

\begin{enumerate}[label=\textbf{\arabic*.}, leftmargin=*, itemsep=0.28em]
  \item Une fonction continue sur un intervalle y admet une unique primitive.
  \item $\displaystyle\int_{a}^{b} f(x)\dx$ est toujours positive.
  \item Si $f \leq g$ sur $\intervalle{a}{b}$ avec $a<b$, alors
        $\int_{a}^{b} f \leq \int_{a}^{b} g$.
  \item La valeur moyenne d'une fonction sur un intervalle est comprise entre
        son minimum et son maximum sur cet intervalle.
\end{enumerate}
\end{jemeteste}

\begin{corrige}[je me teste]
\textbf{1.} Faux — elle en admet une infinité, qui diffèrent d'une constante.
\textbf{2.} Faux — cela dépend du signe de $f$ et de l'ordre des bornes.
\textbf{3.} Vrai.
\textbf{4.} Vrai — conséquence de l'encadrement $m(b-a) \leq \int f \leq M(b-a)$.
\end{corrige}

% ===========================================================================
\begin{commeaubac}
\bmtsource{Exercice au format de l'épreuve — série OSE}
Soit $f$ la fonction définie sur $\intervalle{0}{2}$ par
$f(x) = x\sqrt{4-x^{2}}$.

\begin{enumerate}[label=\textbf{\arabic*.}, leftmargin=*, itemsep=0.25em]
  \item Vérifier que $f$ est positive sur $\intervalle{0}{2}$ et calculer
        $f(0)$ et $f(2)$. \hfill\textup{(0,5 pt)}
  \item Calculer $f'(x)$ et dresser le tableau de variation de $f$ sur
        $\intervalle{0}{2}$. \hfill\textup{(1,5 pt)}
  \item Déterminer une primitive de $f$ sur $\intervalle{0}{2}$.
        \emph{On reconnaîtra une forme $u'u^{n}$.} \hfill\textup{(1 pt)}
  \item En déduire l'aire, en unités d'aire, du domaine délimité par la courbe
        de $f$, l'axe des abscisses et les droites $x=0$ et $x=2$.
        \hfill\textup{(1 pt)}
  \item Calculer la valeur moyenne de $f$ sur $\intervalle{0}{2}$.
        \hfill\textup{(1 pt)}
\end{enumerate}
\end{commeaubac}

\begin{corrige}[exercice au format de l'épreuve]
\textbf{1.} Sur $\intervalle{0}{2}$, $x \geq 0$ et $4-x^{2} \geq 0$ : $f \geq 0$.
$f(0) = f(2) = 0$.

\textbf{2.} $f'(x) = \dfrac{4-2x^{2}}{\sqrt{4-x^{2}}}$, du signe de $2-x^{2}$ :
$f$ croît sur $\intervalle{0}{\sqrt2}$ puis décroît sur $\intervalle{\sqrt2}{2}$,
avec un maximum $f(\sqrt2) = 2$.

\textbf{3.} $f(x) = -\dfrac{1}{2}\times(-2x)\left(4-x^{2}\right)^{1/2}$, forme
$-\frac12 u'u^{1/2}$ avec $u = 4-x^{2}$ : une primitive est
$-\dfrac{1}{3}\left(4-x^{2}\right)^{3/2}$.

\textbf{4.} $\left[-\dfrac{(4-x^{2})^{3/2}}{3}\right]_{0}^{2} = 0 + \dfrac{8}{3}
= \dfrac{8}{3}$ unités d'aire.

\textbf{5.} $\mu = \dfrac{1}{2}\times\dfrac{8}{3} = \dfrac{4}{3}$.
\end{corrige}
